\section{Safety in least sphere degree five}\label{sec:quintic}

Write a sphere phase in the normal monomials \(x^i t^j\) with \(i\ge0\)
and \(y^i t^j\) with \(i>0\). Its leading symbol of degree \(d\), restricted
to the cone \(t^2+2xy=0\), gives a binary form \(h_d(u,v)\) of degree \(2d\)
by \((x,y,t)=(u^2,-v^2/2,uv)\). This identification sends the degree-\(d\)
normal monomials to distinct binary monomials. The \(\mathrm{SL}_2\) action
on \((u,v)\) induces linear isometries of the sphere, by congruence on
\(\left(\begin{smallmatrix}x&t\\t&-2y\end{smallmatrix}\right)\).
They preserve \(\Lambda\), since the sphere trace in each even homogeneous
degree is proportional to the Gaussian trace with covariance inverse to
\(t^2+2xy\), and the odd-degree traces are zero. This proportionality
follows by integrating each normal monomial. In particular we may move a
geometric root of \(h_d\) to either coordinate point after scalar extension.

\subsection{The leading root threshold}\label{sec:leading-threshold}

Suppose every root of \(h_d\) has multiplicity less than \(d\). Embed the
finitely generated coefficient field in \(\mathbb C\). Compactify the sphere
by the smooth projective quadric with boundary conic \(D\). Its affine
residue volume \(\omega\) has divisor \(-2D\) by adjunction. In local
coordinates \((u,z)\) with \(D=\{u=0\}\), write the phase as \(F/u^d\).
At a boundary root of multiplicity \(r>0\), Weierstrass preparation writes
\(F\) as a unit times a monic degree-\(r\) polynomial in \(z\). For fixed
\(u\), the set \(|F|\le\delta\) lies in the union of \(r\) planar disks of
radius \(O(\delta^{1/r})\). Thus on \(|f|\le M\) the squared volume density
has local integral bounded by a constant times
\[
 \int_{|u|<\epsilon}|u|^{-4+2d/r}\,d^2u,
\]
which is finite since \(r<d\). Away from the roots, the region \(|f|\le M\)
stays away from \(D\).

Resolve the pencil and remove its fibre over infinity. On the resulting
smooth surface \(Y\), \(f:Y\to\mathbb A^1\) is proper. The pulled-back
volume form is algebraic and holomorphic on \(Y\), since a divisorial pole
would give an infinite squared integral in a bounded-\(f\) chart.
Take the compact cycle \(\Gamma=\{y=\overline x,\ t\in\mathbb R,\
2|x|^2+t^2=1\}\), which lies in the original affine chart. For
\(-1<t<1\), write \(x=\sqrt{(1-t^2)/2}\,e^{i\theta}\). In these coordinates
the residue form is a nonzero constant multiple of \(d\theta\wedge dt\),
using its expression \(\pm dx\wedge dt/(2x)\) in the chart \(x\ne0\).
Consequently \(\int_\Gamma\omega\ne0\), and integrating \(f^m\omega\)
against \(\Gamma\) gives \((\int_\Gamma\omega)\Lambda(f^m)\).
Since the exponential series converges uniformly on the compact cycle on
bounded \(\tau\)-sets, purity would give the nonzero constant period
\(\int_\Gamma e^{-\tau f}\omega=\int_\Gamma\omega\).
Lemma~\ref{lem:proper-volume}, proved below, excludes this. Thus a pure
phase has a leading root of multiplicity at least \(d\).

Here is the prime-isolation rule used in the remaining cases. Work over
a number field at a good odd prime \(p\) where all coefficients are
integral. If the angular constant term of \(h^m\) is a polynomial
\(C_m(t)\) of degree at most \(p-1\), then
\begin{equation}\label{eq:prime-isolation}
                 p\Lambda(h^m)\equiv [t^{p-1}]C_m(t)\pmod{\mathfrak p}.
\end{equation}
Indeed the integral kills odd heights and divides even-height coefficients
by \(j+1\). The only denominator divisible by \(p\) in the given range is
the denominator at the even height \(p-1\).

An exact multiplicity \(d\) is also impossible under purity. Move that
root to \(u=0\). In the literal Laurent substitution \(y=(1-t^2)/(2x)\),
every monomial \(x^i t^j\) has \(i+j\le d\), and the part on equality is
\(t^d L(x/t)\) for \(L(z)=z^{-d}h_d(z,1)\). With the chosen root, \(L\) is
a polynomial with nonzero constant. At a good prime \(p=1\bmod d\) put
\(m=(p-1)/d\). The coefficient at height \(p-1\) is
\(\operatorname{CT}L^m=L(0)^m\), a unit, so \eqref{eq:prime-isolation}
applies. This argument may be made after
one number-field specialization preserving the leading identity and its
specified nonzero constants. Indeed a finitely generated coefficient
algebra over \(\mathbb Q\), localized to keep these constants nonzero, has a
maximal-ideal quotient that is a number field, and all rational moment
identities specialize. We use the same specialization argument in the
prime tests below, keeping the required distinct roots and support
coefficients nonzero. Primes in the progression \(1\) modulo a positive
integer are infinite, for example by the elementary cyclotomic argument.

\subsection{Supporting faces for multiplicity at least six}\label{sec:face-support}

Let \(d=5\) and move a root of multiplicity at least six to \(v=0\).
All degree-five normal terms now have negative angular exponent, since
their \(v\)-exponents are \(5-i\) for angular exponent \(i\). If the phase
has no positive angular terms, purity forces its zero-angle core to
vanish by the interval moment theorem \cite[Theorem 2.1]{muger-tuset}.
The negative support then gives eventual mixed vanishing for every fixed
marker. Otherwise take the upper supporting line \(j+\rho i=H\) of the
literal Laurent support over angular exponent zero. A convex combination
of nodes reaching angle zero has height less than five, since all
equality nodes of \(i+j\le5\) have negative angle. This gives \(H<5\).
For any degree-five node \(j+\rho i=5+(\rho-1)i\le H\) with \(i<0\),
so \(\rho>1\). In negative angle \(-b\), only the highest Laurent term
of each \(y^b t^r\) can lie on the upper boundary. Thus the possible face
nodes satisfy
\[
 -5\le i\le4,\qquad \max(0,-2i)\le j\le
 \begin{cases}5-i,&i<0,\\4-i,&i\ge0.\end{cases}
\]
A face with no opposite signs must contain a core node at angle zero.
Its height \(H\) is a positive integer, since the support contains nodes
of both signs and all negative-angle nodes have positive height.
At \(m=(p-1)/H\) for good primes \(p\equiv1\bmod H\), the core coefficient
to the \(m\)th power occupies the top angular-constant slot. Equation
\eqref{eq:prime-isolation} excludes it. For a two-node crossing face, write its
angles \(-b,a\) and heights \(j_-,j_+\). Set \(M=aj_-+bj_+\) and choose
a good prime \(p\equiv1\bmod M\), writing \(k=(p-1)/M\).
At exponent \(m=(a+b)k\), angular cancellation uses \(ak\) negative and
\(bk\) positive terms, giving height \(p-1\). The coefficient is the
product of \(\binom{(a+b)k}{ak}\) and a unit monomial in the two
coefficients. The factorials have arguments below \(p\) because
\(M\ge2ab\ge a+b\). Equation \eqref{eq:prime-isolation} excludes this case.

The remaining lines are listed below. To obtain the list, take the line
through every pair of allowed nodes of opposite signs, keep those with
\(\rho>1\), \(H<5\), at least one degree-five node on or below the line,
and at least three equality nodes. An actual crossing face occupies both
signs and at least three listed nodes.
\[
\begin{array}{c|c|l@{\qquad}c|c|l}
\rho&H&i&\rho&H&i\\ \hline
5/3&5/3&-5,-2,1&2&2&-3,-2,-1,0,1\\
3/2&5/2&-5,-3,-1,1&3/2&3&-4,-2,0,2\\
2&3&-2,-1,0,1&3&3&-1,0,1\\
4/3&10/3&-5,-2,1&3/2&7/2&-3,-1,1\\
4/3&11/3&-4,-1,2&5/4&15/4&-5,-1,3\\
4/3&4&-3,0,3&3/2&4&-2,0,2\\
2&4&-1,0,1,2&3/2&9/2&-1,1,3
\end{array}
\]

At \(\rho=H=2\), normal terms have weight at most two for weights
\((x,y,t)=(2,0,1)\). After exchanging \(x,y\), the literal Laurent height
is at most two. Regularity then allows at most a simple negative angular
pole. Section~\ref{sec:quadratic-height} proves safety for every such pure
phase of least sphere degree five.

\subsection{The cyclic-cover obstruction on a face}\label{sec:cyclic-face}

For a noninteger height \(H=n/q\) in lowest terms, the angular Laurent
polynomial of the face is \(z^{-b}P(z^q)\), where \(-b,a\) are its occupied
extreme angles. The table gives the same \(q\) as the denominator of \(\rho\).
In an integer-height row, first divide the angles by their common step if
the displayed angles are all multiples of two or three, and then put
\(q=1\) and \(n=H\). Angular compression leaves constant terms unchanged.
In these coordinates \(\deg P=(a+b)/q\), because both extreme terms
are occupied. For a prime \(p=nk+1\) and moment order \(qk\), the constant angular term has
height at most \(qkH=p-1\). A product contributing at that height uses
only equality nodes of the supporting line. The only interval coefficient
with denominator divisible by \(p\) is the one of height \(p-1\). Thus
multiplying the moment by \(p\) and reducing gives
\[
 [z^{bk}]P(z)^{qk}=[z^{p-1}]\bigl(z^{n-b}P(z)^q\bigr)^k.
\tag{*}\label{eq:face-scalar}
\]
The raw degree-\(2n\) branch form in the second expression has endpoint
exponents \(n-b,n-a\), and each root of \(P\) has its multiplicity multiplied
by \(q\). When all nonzero exponents are less than \(n\), the normalized
cyclic cover \(w^n=z^{n-b}P(z)^q\) detects \eqref{eq:face-scalar} in its
character-one \(H^1(O)\) line. In fact the character-one sheaf is \(O(-2)\)
and its \v{C}ech generator \(w/z\) maps to
\(wz^{-p}(z^{n-b}P(z)^q)^k\) under Frobenius. Projection to the line gives
exactly the selected coefficient. More generally the character-\(i\)
dimension is \(-1+\sum_j\{i a_j/n\}\) for branch exponents \(a_j\), by
dividing the local Kummer generator \(w^i\) by the branch coordinate to
power \(\lfloor i a_j/n\rfloor\). All the covers used here are connected,
as the branch exponents have greatest common divisor one with \(n\).
The rational idempotents of the cyclic deck group give Jacobian factors
up to isogeny. The factor for a divisor \(s>1\) of \(n\) has action by
\(\mathbb Q(\zeta_s)\), and its dimension is the sum of the coherent
character dimensions for characters of order \(s\). Complex conjugation
interchanges the coherent and holomorphic-one-form characters, preserving
the unordered multiplicities of each conjugate pair.
Ordinary reduction makes Frobenius invertible on coherent \(H^1\).
It suffices for the primitive factor to be ordinary because the
character-one line belongs to that factor. The nonzero scalar then
excludes the face by \eqref{eq:prime-isolation} and \eqref{eq:face-scalar}.
Work away from the finitely many primes of bad reduction for the curve,
cover, and deck action. Adjoin the required roots of unity and take
ordinary primes of absolute degree one, so they lie in the required
prime progression.

For heights three and four the compressed angular sets, up to sign, are
subsets of \(\{-2,-1,0,1\}\) or \(\{-1,0,1\}\). The branch exponents
are displayed below. The first two cases of the cubic partition are
all distinct and one double root. The third is a triple root.
\[
\begin{array}{c|c|c|c}
H & \deg P & \text{distinct roots} & \text{coincident roots}\\ \hline
3 & 2 &(2,1,1,2)&(2,2,2)\\
3 & 3 &(1,1,1,1,2)&(1,2,1,2),\ (1,3,2)\\
4 & 2 &(3,1,1,3)&(3,2,3)\\
4 & 3 &(2,1,1,1,3)&(2,2,1,3),\ (2,3,3)
\end{array}
\]
At height three the distinct quadratic and cubic root configurations
give genus two and three, respectively. In genus three the Jacobian has imaginary quadratic action
by \(\mathbb Q(\zeta_3)\). At height four the distinct-root quadratic
configuration has genus three, with an elliptic order-two factor and a
surface with order-four action. The distinct-root cubic configuration
has genus four, with an elliptic order-two factor and an imaginary
quadratic threefold. All repeated configurations have genus at most two,
except the triple cubic root at height three whose raw exponent equals
three. There \eqref{eq:face-scalar} is a unit times \(\binom{3k}{2k}\).
The genus at most two cases have ordinary primes of density one after
finite extension \cite{sawin}. For the genus-four case, first pass to a
common extension containing the roots of unity over which the elliptic
factor has density-one ordinarity. Apply Fit\'e's theorem over that same
field to the threefold with imaginary quadratic multiplication.
Its positive-density ordinary set intersects both the elliptic ordinary
set and the absolute degree-one primes in a set of positive density.
The same theorem gives the required primes for the genus-three case
\cite[Theorem~1]{fite}.

The noninteger heights with \(n\in\{7,9,10,11,15\}\) have quadratic \(P\).
For distinct roots the primitive signatures in increasing character order
are as follows.
\[
\begin{array}{c|c|c}
n&\text{branch exponents}&\text{primitive signature}\\ \hline
7 &(4,2,2,6)&(1,1,2,0,1,1)\\
9 &(8,2,2,6)&(1,1,2,0,1,1)\\
10&(5,3,3,9)&(1,2,0,1)\\
11&(7,3,3,9)&(1,1,2,0,0,2,2,0,1,1)\\
15&(10,4,4,12)&(1,1,0,2,0,2,1,1)
\end{array}
\]
In each row the primitive factor has rank two over \(\mathbb Q(\zeta_n)\)
and at most two conjugate pairs of multiplicities \((1,1)\), so
Theorem~\ref{lem:rank-two} applies. The deck action preserves the
Jacobian polarization, so its Rosati involution inverts the deck generator
and induces complex conjugation on \(\mathbb Q(\zeta_n)\). For \(n=9\)
there is also an elliptic factor of order three. For \(n=10\) the
additional factors have orders two and five, with dimensions one and
two, respectively. For \(n=15\) they have orders three and five, again
with dimensions one and two. There are no lower-order factors for
\(n=7,11\). If the roots coincide, merge the two equal interior
exponents. The character formula gives dimension zero at every lower
order in all five cases. The primitive factor has rank one over its
cyclotomic field. After
adjoining the field, its crystalline components at good degree-one primes
have height one and are ordinary. This proves exclusion in all five
rows, including the repeated-root cases.

For example, the height \(7/2\) face has angles \(-3,-1,1\).
Writing it as \(z^{-3}P(z^2)\), at \(p=7k+1\) the moment of order \(2k\)
detects \([z^{3k}]P(z)^{2k}\). The associated cover is
\(w^7=z^4P(z)^2\). When \(P\) has distinct roots, the Jacobian is a
sixfold with rank-two \(\mathbb Q(\zeta_7)\)-action. Its character pairs
\((1,6)\) and \((2,5)\) are balanced, while \((3,4)\) has multiplicities
\((2,0)\). The two balanced pairs are precisely the situation in which
the extra trace divisibility in Theorem~\ref{lem:rank-two} is used.

\subsection{The order-five collisions}

At height \(5/3\) a distinct-root quadratic gives exponents \((3,3,4)\) and
genus two. At height \(5/2\) the positive extreme is \(1\) and the negative
extreme is \(-3\) or \(-5\). In both cases distinct roots give exponents
\((4,2,2,2)\) and a rank-two order-five Jacobian, so
Theorem~\ref{lem:rank-two} applies. Merging two of the exponents \(2\)
gives genus two. There remain only the raw exponents exceeding five:
the repeated quadratic root at height \(5/3\) and the triple cubic root
at height \(5/2\).

For these cases use the sphere automorphism
\[
 \theta_{c,r}:(x,y,t)\longmapsto
 \bigl(x-cy^r t-c^2y^{2r+1}/2,\ y,\ t+cy^{r+1}\bigr).
\]
It preserves \(\Lambda\): its locally nilpotent derivation in the
Laurent \(y\)-chart is \(cy^{r+1}\partial_t\), and integration gives
zero by the endpoint formula, since \(x=0\) at \(t=\pm1\). Thus it also
preserves purity and safety by transporting markers.

At height \(5/3\), the face is \(Bx+Ay^2t+Cy^5\), with
\(A^2+2BC=0\) in the repeated-root case. It is the leading part of
\(B\theta_{-A/B,2}(x)\) for weights \((5,1,3)\). The inverse shear gives
\(Bx+P_4(y)+(a+b y)t\), where \(\deg P_4\le4\). The possible \(y^4\) term
is excluded by the two-node test. The next possible tie is between \(y^3\)
and \(yt\), since both have slope \(3/2\) from the node of \(Bx\). If both
occur, a distinct-root face gives the order-three cover with exponents
\((2,2,2)\), and a repeated-root face is the leading part of
\(B\theta_{-b/B,1}(x)\) for weights \((3,1,2)\). Invert this shear as well.
Any pure phase left has degree at most two.

At height \(5/2\), the triple-root face is the leading part of
\(B\theta_{c,1}(xt)\) for weights \((3,1,2)\). Indeed its face cubic is
\(B(z+c/4)^3\), as follows by expanding the shear and replacing
\(xy=(1-t^2)/2\). The inverse shear gives
\(x(Bt+a)+P_4(y)+tP_2(y)+d t^2\) with \(\deg P_i\le i\). The slopes from
\(Bxt\) to \(y^4,y^2t,y^3\) are \(7/5,4/3,5/4\), respectively, each strictly
largest once the earlier terms have been removed. The two-node test
removes all three. Again the remaining degree is at most two.

Section~\ref{sec:quadratic-height} gives safety for every pure phase of
degree at most two. This completes the unstable-leading branch and proves
safety for every pure phase of least sphere degree five.
